\documentclass[10pt,a4paper]{amsart}
\usepackage[margin=19mm]{geometry}
\usepackage{amsmath,amssymb,mathtools}
\usepackage[scr=rsfs,cal=euler]{mathalfa}
\usepackage{xfrac}
\usepackage[unicode,hidelinks,bookmarks=false]{hyperref}
\newcommand{\ie}{i.e.\ }
\newcommand{\cf}{cf.\ }
\newcommand{\resp}{respectively\ }
\newcommand{\Subsection}{\subsection}
\providecommand{\nobreakdash}{\nobreak\hbox{-}}
\setlength{\emergencystretch}{2em}
\multlinegap0pt
\usepackage{bm}
\usepackage{bbm}
\usepackage{dsfont}
\usepackage{xspace}
\usepackage{cancel}
\usepackage{mathtools}
\usepackage[shortlabels]{enumitem}
\usepackage{amsthm}
\makeatletter
\@ifundefined{theorem}{\newtheorem{theorem}{Theorem}}{}
\@ifundefined{lemma}{\newtheorem{lemma}[theorem]{Lemma}}{}
\makeatother
\newcommand{\acc}{\mathrm{acc}}
\newcommand{\mc}{\mathcal}
\newcommand{\nacc}{\mathrm{nacc}}
\title[Higher walks and squares]{Higher walks and squares}
\author{Chris Lambie-Hanson, Pedro Marun}
\date{Source: arXiv:2512.08633v2 (CC BY 4.0)}
\begin{document}
\maketitle

\noindent\emph{Source and licence.} Higher walks and squares. Version arXiv:2512.08633v2 (CC BY 4.0), distributed under the Creative Commons Attribution 4.0 International licence, as recorded in the arXiv licence metadata for this exact version and shown on the abstract page https://arxiv.org/abs/2512.08633v2. Full rights notice: licenses/arxiv-2512.08633-CC-BY-4.0.txt.\par\smallskip

\noindent\textit{Editorial scope.} Counted unit: the Lemma whose source label is \texttt{simulation} in the article (source0.tex lines 1983--2000, source label \texttt{simulation}), delivered together with the article's own complete original proof of it (source0.tex lines 2002--2135, one \texttt{proof} environment of 134 lines). No mathematics is written, repaired or re-derived: statement and proof are the article's own text line for line. Required context retained uncounted: none. Every definition, display and label of those lines is retained unchanged. Omitted from this package and not redistributed: 1--1982, 2001--2001, 2136--2596. Nothing inside a retained excerpt is omitted or altered: every retained body is delivered exactly as the article's lines are written, so no omission receipt of any kind accompanies this package. Citations: the retained excerpts carry no citation command at all, so no bibliography entry of the article is transcribed and no reference is redistributed. Reference adaptation: none was needed and none was made; every reference inside the delivered unit resolves inside it, so no unresolved reference or citation is produced. Printed slips kept literal: no printed word, symbol, hypothesis or number of the counted statement or of its proof was altered. Layout and class: the article's own class is \texttt{amsart} with packages including paralist, color, mathrsfs, amssymb, amsmath, amsfonts, stmaryrd, graphics, enumitem, mathdots, showlabels, float, changes, phoenician; the wrapper is the lane's pinned \texttt{amsart} editorial wrapper at 10pt on A4, which restates the theorem-like environments the retained text needs and adds the article's own macro definitions that the retained text needs (\texttt{\textbackslash acc}, \texttt{\textbackslash mc}, \texttt{\textbackslash nacc}), with extra wrapper packages (bm, bbm, dsfont, xspace, cancel, mathtools, enumitem). The delivered numbering is the unit's own. Assets: no retained span contains a figure environment, a graphics-inclusion command, a PDF-page inclusion, a table or a TikZ picture; no image file is copied into this package, the package declares no asset dependency and no image file is redistributed with it. Licence: version arXiv:2512.08633v2 of this article is distributed under the Creative Commons Attribution 4.0 International licence, as recorded in the arXiv licence metadata for this exact version and shown on the abstract page https://arxiv.org/abs/2512.08633v2. A full rights notice is delivered at licenses/arxiv-2512.08633-CC-BY-4.0.txt. Characters: every retained excerpt is pure ASCII, so the delivered engine, XeLaTeX with its default Latin Modern text font, reports no missing character.
\begin{lemma}\label{simulation}
  Suppose that $m$ is a positive integer and $\mc{C}$ is a coherent 
  $(m+2)$-sequence on an ordinal $\lambda$. Suppose that $\vec{\gamma} \in 
  \lambda^{[m]}$ and $\alpha \leq \gamma_0$ is in $X(\mc{C})$.
  Suppose moreover that $L_m^-(\alpha,\vec{\gamma}) < \xi < \alpha$, and
  let $\eta_\xi = \min(C_\alpha \setminus \xi)$. Then
  \begin{enumerate}
    \item $S_{m+2}(\xi,\eta_\xi,\alpha,\vec{\gamma}) = 
    (S^-_m(\alpha,\vec{\gamma}))^{m+2}$; and
    \item for all $x \in S^-_m(\alpha,\vec{\gamma})$, if 
    $\mathrm{Tr}_m(+,\alpha,\vec{\gamma})(x) = ((-1)^j,\alpha,\vec{\beta})$, 
    then 
    \[
      \mathrm{Tr}_{m+2}(+,\xi,\eta_\xi,\alpha,\vec{\gamma})(s_{m,m+2}(x)) = 
      ((-1)^j,\xi,\eta_\xi,\alpha,\vec{\beta}).
    \]
  \end{enumerate}
\end{lemma}

\begin{proof}
  We will prove the following statements for all $x \in S^-_m(\alpha,\vec{\gamma})$ 
  by induction on $|x|$:
  \begin{enumerate}[label = (\alph*)]
    \item $s_{m,m+2}(x) \in S_{m+2}(\xi,\eta_\xi,\alpha,\vec{\gamma})$;
    \item if $\mathrm{Tr}_m(+,\alpha,\vec{\gamma})(x) = 
    ((-1)^j,\alpha,\vec{\beta})$, then 
    $\mathrm{Tr}_{m+2}(+,\xi,\eta_\xi,\alpha,\vec{\gamma})(s_{m,m+2}(x)) = 
    ((-1)^j,\xi,\eta_\xi,\alpha,\vec{\beta})$;
    \item $s_{m,m+2}(x)$ is a terminal node of $S_{m+2}(\xi,\eta_\xi,\alpha, 
    \vec{\gamma})$ if and only if $x$ is a terminal node of 
    $S^-_m(\alpha,\vec{\gamma})$;
    \item if $x$ is a splitting node of $S^-_m(\alpha,\vec{\gamma})$, then, for $i < 2$, 
    $s_{m,m+2}(x)^\smallfrown \langle i \rangle$ is a terminal node of 
    $S_{m+2}(\xi,\eta_\xi,\alpha,\vec{\gamma})$.
  \end{enumerate}
  This will clearly suffice to establish the lemma.
  We begin with the base case, i.e., $x = \emptyset$. Items (a) and (b) are 
  immediate. To see (c), let $\tau(\vec{\gamma})$ denote the longest final segment 
  of $\vec{\gamma}$ that is in $I(\mc{C})$. Let 
  $k = |\tau(\vec{\gamma})|$ and $j = (m-1)-k$. In particular, setting 
  $\alpha = \gamma_{-1}$ for convenience, $\gamma_j$ is the largest ordinal appearing 
  in $\langle \alpha \rangle^\smallfrown \vec{\gamma}$ that does not appear in 
  $\tau(\vec{\gamma})$. There are now a number of options:
  \begin{itemize}
    \item If $\tau(\vec{\gamma}) = \vec{\gamma}$ and $\alpha \in C_{\vec{\gamma}}$, 
    then $x$ is a terminal node of $S^-_m(\alpha,\vec{\gamma})$. There are now two 
    subcases to consider:
    \begin{itemize}
      \item If $\alpha \in \acc(C_{\vec{\gamma}})$, then we have 
      $C_{\alpha\vec{\gamma}}= C_\alpha$. In particular, 
      $\langle \eta_\xi,\alpha \rangle ^\smallfrown \vec{\gamma} \in I(\mc{C})$ 
      and $C_{\eta_\xi\alpha\vec{\gamma}} \setminus \xi = \emptyset$, so $x = s_{m,m+2}(x)$ 
      is a terminal node of $S_{m+2}(\xi,\eta_\xi,\alpha,\vec{\gamma})$.
      \item If $\alpha \in \nacc(C_{\vec{\gamma}})$, then we have 
      $\max(C_{\vec{\gamma}} \cap \alpha) \leq L^-_m(\alpha,\vec{\gamma}) < 
      \xi < \eta_\xi < \alpha$. In particular, $\langle \alpha \rangle ^\smallfrown 
      \vec{\gamma}$ is the longest final segment of $\langle \xi, \eta_\xi, 
      \alpha, \vec{\gamma} \rangle$ in $I(\mc{C})$ and $C_{\alpha\vec{\gamma}} 
      \setminus \eta_\xi = \emptyset$, so again $x$ is a terminal node of 
      $S_{m+2}(\xi,\eta_\xi,\alpha,\vec{\gamma})$.
    \end{itemize}
    \item If $C_{\tau(\vec{\gamma})} \setminus \gamma_j = \emptyset$, then again 
    $x$ is a terminal node of $S^-_m(\alpha,\vec{\gamma})$. In this case, 
    $\tau(\vec{\gamma})$ is also the longest final segment of 
    $\langle \xi,\eta_\xi,\alpha \rangle^\smallfrown \vec{\gamma}$ that is in 
    $I(\mc{C})$ and $\gamma_j$ is the largest element of $\langle \xi, 
    \eta_\xi, \alpha \rangle ^\smallfrown \vec{\gamma}$ not appearing in 
    $\tau(\vec{\gamma})$, so it immediately follows 
    that $x$ is also a terminal node of $S_{m+2}(\xi,\eta_\xi,\alpha,\vec{\gamma})$.
    \item If we are in none of the above cases, then $x$ is a splitting node 
    of $S^-_m(\alpha,\vec{\gamma})$. In particular, we have $\gamma_j \notin 
    C_{\tau(\vec{\gamma})}$ and $C_{\tau(\vec{\gamma})} \setminus \gamma_j 
    \neq \emptyset$. As in the previous item, we again see that $\tau(\vec{\gamma})$ 
    is the longest final segment of $\langle \xi,\eta_\xi,\alpha \rangle^\smallfrown 
    \vec{\gamma}$ that is in $I(\mc{C})$, and the fact that 
    $C_{\tau(\vec{\gamma})} \setminus \gamma_j \neq \emptyset$ implies that 
    $x$ is a splitting node of $S_{m+2}(\xi,\eta_\xi,\alpha,\vec{\gamma})$.
  \end{itemize}
  We finally verify (d). If $x$ is a splitting node of $S^-_m(\alpha, \vec{\gamma})$, 
  then we are in the case of the final bullet point above. Let 
  $\beta = \min(C_{\tau(\vec{\gamma})} \setminus \gamma_j)$, and note that 
  $\gamma_j < \beta$. In this case, recalling the definition of $\mathrm{Tr}_{m+2}$, we have
  \[
    \mathrm{Tr}_{m+2}(\xi,\eta_\xi,\alpha,\vec{\gamma})(\langle 0 \rangle) = 
    (\xi, \alpha, \gamma_0, \ldots, \gamma_j, \beta, \tau(\vec{\gamma})),
  \]
  which is terminal due to the fact that $\langle \beta \rangle ^\smallfrown 
  \tau(\vec{\gamma}) \in I(\mc{C})$ and $C_{\beta \tau(\vec{\gamma})} \setminus 
  \gamma_j = \emptyset$.
  
  If $j > -1$, then we have
  \[
    \mathrm{Tr}_{m+2}(\xi,\eta_\xi,\alpha,\vec{\gamma})(\langle 1 \rangle) = 
    (\xi, \eta_\xi, \gamma_0, \ldots, \gamma_j, \beta, \tau(\vec{\gamma})),
  \]
  which is terminal for the same reason. If $j = -1$, and hence 
  $\tau(\vec{\gamma}) = \vec{\gamma}$, then we have
  \[
    \mathrm{Tr}_{m+2}(\xi,\eta_\xi,\alpha,\vec{\gamma})(\langle 1 \rangle) = 
    (\xi, \eta_\xi, \beta, \tau(\vec{\gamma})).
  \]
  In this case, we know that $\max(C_{\tau(\vec{\gamma})} \cap \alpha) 
  \leq L^-_m(\alpha,\vec{\gamma}) < \xi < \eta_\xi < \alpha$, and hence we have 
  $\langle \beta \rangle ^\smallfrown \tau(\vec{\gamma}) \in I(\mc{C})$ and 
  $C_{\beta \tau(\vec{\gamma})} \setminus \eta_\xi = \emptyset$, so again 
  $\langle 1 \rangle$ is terminal in $S_{m+2}(\xi,\eta_\xi,\alpha,\vec{\gamma})$. 
  This completes the base case of the induction.
  
  Now suppose that $x \in S^-_m(\alpha, \vec{\gamma})$ is a splitting node and 
  we have established (a)--(d) for $x$. Fix $i < m$; we will establish 
  (a)--(d) for $x^\smallfrown \langle i \rangle$. First note that 
  $s_{m,m+2}(x^\smallfrown \langle i \rangle) = s_{m,m+2}(x)^\smallfrown \langle i+2 \rangle$. 
  Item (a) now follows immediately from the fact that, by the inductive hypothesis, 
  $s_{m,m+2}(x)$ is a splitting node of $S_{m+2}(\xi,\eta_\xi,\alpha,\vec{\gamma})$.
  
  To see (b), let $\mathrm{Tr}_m(+,\alpha,\vec{\gamma})(x) = ((-1)^\ell, \alpha, 
  \vec{\beta})$. As in the base case, let $\tau(\vec{\beta})$ be the longest final 
  segment of $\vec{\beta}$ that is in $I(\mc{C})$, let $k = |\tau(\vec{\beta})|$, 
  let $j = (m-1)-k$ and, for convenience, set $\beta_{-1} = \alpha$. Since $x$ is a 
  splitting node, it must be the case that $\beta_j \notin C_{\tau(\vec{\beta})}$ 
  and $C_{\tau(\vec{\beta})} \setminus \beta_j \neq \emptyset$. By the inductive 
  hypothesis, we have $\mathrm{Tr}_{m+2}(+,\xi,\eta_\xi,\alpha,\vec{\gamma})
  (s_{m,m+2}(x)) = ((-1)^\ell, \xi, \eta_\xi, \alpha, \vec{\beta})$; moreover, $\tau(\vec{\beta})$ 
  must be the longest final segment of $\langle \xi,\eta_\xi,\alpha \rangle^\frown \vec{\beta}$ 
  that is in $I(\mc{C})$.
  
  Let $\delta = \min(C_{\tau(\vec{\beta})} \setminus \beta_j)$. If $i \leq j$, then 
  \[
    \mathrm{Tr}_m(+,\alpha,\vec{\gamma})(x^\smallfrown \langle i \rangle) = 
    ((-1)^{\ell+j+i},\alpha, (\vec{\beta} \restriction (j+1))^i, \delta, \tau(\vec{\beta}))
  \]
  and 
  \[
    \mathrm{Tr}_{m+2}(+,\xi,\eta_\xi,\alpha,\vec{\gamma})(s_{m,m+2}(x)^\smallfrown 
    \langle i+2 \rangle) = ((-1)^{\ell+j+i}, \xi,\eta_\xi,\alpha, (\vec{\beta} \restriction (j+1))^i, 
    \delta, \tau(\vec{\beta})).
  \]
  If $i > j$, then 
  \[
    \mathrm{Tr}_m(+,\alpha,\vec{\gamma})(x^\smallfrown \langle i \rangle) = 
    ((-1)^{\ell+j+i+1}, \alpha, \vec{\beta} \restriction (j+1), \delta, \tau(\vec{\beta})^{i-(j+1)})
  \]
  and 
  \[
    \mathrm{Tr}_{m+2}(+,\xi,\eta_\xi,\alpha,\vec{\gamma})(s_{m,m+2}(x)^\smallfrown 
    \langle i+2 \rangle) = ((-1)^{\ell+j+i+1}, \xi,\eta_\xi,\alpha, \vec{\beta} \restriction (j+1), 
    \delta, \tau(\vec{\beta})^{i-(j+1)}).
  \]
  In either case, $x^\smallfrown \langle i \rangle$ satisfies item (b).
  
  Items (c) and (d) are now verified exactly as in the base case, \emph{mutatis 
  mutandis}. We therefore leave this to the reader.
\end{proof}

\end{document}
